% only the text for the summary

\paragraph{Motivation.}
Kubernetes autoscaling has become a fundamental operational challenge as cloud-native
applications grow in scale and complexity. The two dominant mechanisms—Horizontal
Pod Autoscaler (HPA) and Vertical Pod Autoscaler (VPA)—operate on fixed CPU/memory
thresholds and reactive control loops. They act after load spikes arrive, often too
late to prevent SLA violations. They over-provision resources to hedge against future
bursts. Most critically, they produce no explanation for their decisions: an operator
observing an unexpected scale-up event has no record of the reasoning behind it.
This opacity is increasingly problematic in regulated environments and in
organizations seeking to audit their infrastructure cost.

\paragraph{Contributions.}
This thesis introduces \textbf{Ai4k8s}, an open-source agentic AI platform hosted
at \url{https://ai4k8s.online}, comprising a web-based Kubernetes management interface,
and an MCP (Model Context Protocol) tool server for Kubernetes operations.
The principal contributions are:

\begin{enumerate}
  \item \textbf{A three-layer decision pipeline.} Ai4k8s combines a safety enforcement
  layer, a local LLM reasoning layer, and an MCDA optimization layer (TOPSIS). Every
  scaling decision is the product of all three layers and is logged with full provenance:
  which layer fired, what the LLM said, and what MCDA scores were assigned to each
  candidate action.

  \item \textbf{On-premises LLM autoscaling.} We demonstrate that a 522~MB
  quantized language model (\textit{qwen3:0.6b}) running on a single-node cluster
  with no GPU is sufficient to produce valid, structured autoscaling recommendations
  in 1.7~s at idle. A resilient cascade (local $\to$ cloud LLM $\to$ regex fallback)
  ensures continuous operation under CPU saturation.

  \item \textbf{52\% cost reduction vs.\ HPA.} On a three-service muBench workload
  at 48 concurrent connections, Ai4k8s achieves a normalized cost proxy of 0.184
  compared to 0.380 for native HPA—a 52\% reduction—by avoiding unnecessary
  scale-up actions when the forecast does not justify the added replicas.

  \item \textbf{Explainability.} Ai4k8s is the first Kubernetes autoscaler to
  produce a structured, human-readable decision trace for every scaling action,
  including LLM chain-of-thought reasoning and per-criterion MCDA scores.
\end{enumerate}

\paragraph{Evaluation.}
We evaluate on a CrownLabs VM (4 vCPU, 14~GiB RAM, k3s single-node) running a
muBench three-service chain (\textit{ingest $\to$ process $\to$ analyze}) under
HTTP burst traffic generated by \textit{wrk}. Structured experiments cover bottleneck
analysis (seq\_len tuning), predictive scale-up with forced rising forecasts, HPA/VPA
mode selection accuracy (3/3 correct), and a full N=3 comparison across 48c and 96c
connection levels. The evaluation confirms that Ai4k8s correctly maintains at
sub-threshold CPU (avoiding HPA's reactive over-provisioning) while the LLM and
MCDA layers surface early scale-up signals ahead of HPA reaction time.

\paragraph{Limitations and future work.}
On a 4-vCPU VM the local LLM competes with Kubernetes pods for CPU time under full
load, triggering the Groq cloud fallback. An 8-vCPU instance is projected to allow
local-only inference under production loads. Future work includes multi-cluster
deployment, fine-tuning the LLM on Kubernetes metrics corpora, and integrating
reinforcement learning to adapt TOPSIS weights from observed outcomes.

% $400\times$ is nicer than 400x
